% dataset: math-formula-atlas

% page_slug: exponent-laws-html

% page_title: Laws of Exponents

% page_url: https://www.mathsisfun.com/algebra/exponent-laws.html

% subject_slug: algebra

% subject_path: /algebra/

% formula_count: 42

% formula_record_start
% formula_id: exponent-laws-html-001
% index_global: 1
% index_in_page: 1
% section: Laws of Exponents
% subsection: none
% formula_name: 
% source: mathsisfun-large-span
\[x^{1} = x\]
% formula_record_end

% formula_record_start
% formula_id: exponent-laws-html-002
% index_global: 2
% index_in_page: 2
% section: Laws of Exponents
% subsection: none
% formula_name: 
% source: mathsisfun-large-span
\[x^{0} = 1\]
% formula_record_end

% formula_record_start
% formula_id: exponent-laws-html-003
% index_global: 3
% index_in_page: 3
% section: Laws of Exponents
% subsection: none
% formula_name: 
% source: mathsisfun-large-span
\[x^{-1} = 1/x\]
% formula_record_end

% formula_record_start
% formula_id: exponent-laws-html-004
% index_global: 4
% index_in_page: 4
% section: Laws of Exponents
% subsection: none
% formula_name: 
% source: mathsisfun-large-span
\[x^{m}x^{n} = x^{m+n}\]
% formula_record_end

% formula_record_start
% formula_id: exponent-laws-html-005
% index_global: 5
% index_in_page: 5
% section: Laws of Exponents
% subsection: none
% formula_name: 
% source: mathsisfun-large-span
\[x^{m}/x^{n} = x^{m-n}\]
% formula_record_end

% formula_record_start
% formula_id: exponent-laws-html-006
% index_global: 6
% index_in_page: 6
% section: Laws of Exponents
% subsection: none
% formula_name: 
% source: mathsisfun-large-span
\[(x^{m})^{n} = x^{mn}\]
% formula_record_end

% formula_record_start
% formula_id: exponent-laws-html-007
% index_global: 7
% index_in_page: 7
% section: Laws of Exponents
% subsection: none
% formula_name: 
% source: mathsisfun-large-span
\[(xy)^{n} = x^{n}y^{n}\]
% formula_record_end

% formula_record_start
% formula_id: exponent-laws-html-008
% index_global: 8
% index_in_page: 8
% section: Laws of Exponents
% subsection: none
% formula_name: 
% source: mathsisfun-large-span
\[(x/y)^{n} = x^{n}/y^{n}\]
% formula_record_end

% formula_record_start
% formula_id: exponent-laws-html-009
% index_global: 9
% index_in_page: 9
% section: Laws of Exponents
% subsection: none
% formula_name: 
% source: mathsisfun-large-span
\[x^{-n} = 1/x^{n}\]
% formula_record_end

% formula_record_start
% formula_id: exponent-laws-html-010
% index_global: 10
% index_in_page: 10
% section: Laws of Exponents
% subsection: none
% formula_name: 
% source: mathsisfun-large-span
\[x^{1} = x\]
% formula_record_end

% formula_record_start
% formula_id: exponent-laws-html-011
% index_global: 11
% index_in_page: 11
% section: Laws of Exponents
% subsection: none
% formula_name: 
% source: mathsisfun-large-span
\[x^{0} = 1\]
% formula_record_end

% formula_record_start
% formula_id: exponent-laws-html-012
% index_global: 12
% index_in_page: 12
% section: Laws of Exponents
% subsection: none
% formula_name: 
% source: mathsisfun-large-span
\[x^{-1} = 1/x\]
% formula_record_end

% formula_record_start
% formula_id: exponent-laws-html-013
% index_global: 13
% index_in_page: 13
% section: Laws of Exponents
% subsection: none
% formula_name: x1 = x
% source: mathsisfun-table-law
\textbf{Label: x1 = x}
\[x^{1} = x\]
% formula_record_end

% formula_record_start
% formula_id: exponent-laws-html-014
% index_global: 14
% index_in_page: 14
% section: Laws of Exponents
% subsection: none
% formula_name: Example: x1 = x
% source: mathsisfun-table-example
\textbf{Label: Example: x1 = x}
\[6^{1}=6\]
% formula_record_end

% formula_record_start
% formula_id: exponent-laws-html-015
% index_global: 15
% index_in_page: 15
% section: Laws of Exponents
% subsection: none
% formula_name: x0 = 1
% source: mathsisfun-table-law
\textbf{Label: x0 = 1}
\[x^{0} = 1\]
% formula_record_end

% formula_record_start
% formula_id: exponent-laws-html-016
% index_global: 16
% index_in_page: 16
% section: Laws of Exponents
% subsection: none
% formula_name: Example: x0 = 1
% source: mathsisfun-table-example
\textbf{Label: Example: x0 = 1}
\[7^{0}=1\]
% formula_record_end

% formula_record_start
% formula_id: exponent-laws-html-017
% index_global: 17
% index_in_page: 17
% section: Laws of Exponents
% subsection: none
% formula_name: x-1 = 1/x
% source: mathsisfun-table-law
\textbf{Label: x-1 = 1/x}
\[x^{-1} = 1/x\]
% formula_record_end

% formula_record_start
% formula_id: exponent-laws-html-018
% index_global: 18
% index_in_page: 18
% section: Laws of Exponents
% subsection: none
% formula_name: Example: x-1 = 1/x
% source: mathsisfun-table-example
\textbf{Label: Example: x-1 = 1/x}
\[4^{-1}=\frac{1}{4}\]
% formula_record_end

% formula_record_start
% formula_id: exponent-laws-html-019
% index_global: 19
% index_in_page: 19
% section: Laws of Exponents
% subsection: none
% formula_name: xmxn = xm+n
% source: mathsisfun-table-law
\textbf{Label: xmxn = xm+n}
\[x^{m}x^{n} = xm+n\]
% formula_record_end

% formula_record_start
% formula_id: exponent-laws-html-020
% index_global: 20
% index_in_page: 20
% section: Laws of Exponents
% subsection: none
% formula_name: Example: xmxn = xm+n
% source: mathsisfun-table-example
\textbf{Label: Example: xmxn = xm+n}
\[x^{2}x^{3}=x^{2+3}=x^{5}\]
% formula_record_end

% formula_record_start
% formula_id: exponent-laws-html-021
% index_global: 21
% index_in_page: 21
% section: Laws of Exponents
% subsection: none
% formula_name: xm/xn = xm-n
% source: mathsisfun-table-law
\textbf{Label: xm/xn = xm-n}
\[\frac{x^{m}}{x^{n}} = xm-n\]
% formula_record_end

% formula_record_start
% formula_id: exponent-laws-html-022
% index_global: 22
% index_in_page: 22
% section: Laws of Exponents
% subsection: none
% formula_name: Example: xm/xn = xm-n
% source: mathsisfun-table-example
\textbf{Label: Example: xm/xn = xm-n}
\[x^{6}/x^{2} = x^{6}-2 = x^{4}\]
% formula_record_end

% formula_record_start
% formula_id: exponent-laws-html-023
% index_global: 23
% index_in_page: 23
% section: Laws of Exponents
% subsection: none
% formula_name: (xm)n = xmn
% source: mathsisfun-table-law
\textbf{Label: (xm)n = xmn}
\[(x^{m})^{n} = xmn\]
% formula_record_end

% formula_record_start
% formula_id: exponent-laws-html-024
% index_global: 24
% index_in_page: 24
% section: Laws of Exponents
% subsection: none
% formula_name: Example: (xm)n = xmn
% source: mathsisfun-table-example
\textbf{Label: Example: (xm)n = xmn}
\[(x^{2})3 = x^{2}\times 3 = x^{6}\]
% formula_record_end

% formula_record_start
% formula_id: exponent-laws-html-025
% index_global: 25
% index_in_page: 25
% section: Laws of Exponents
% subsection: none
% formula_name: (xy)n = xnyn
% source: mathsisfun-table-law
\textbf{Label: (xy)n = xnyn}
\[(xy)^{n} = xnyn\]
% formula_record_end

% formula_record_start
% formula_id: exponent-laws-html-026
% index_global: 26
% index_in_page: 26
% section: Laws of Exponents
% subsection: none
% formula_name: Example: (xy)n = xnyn
% source: mathsisfun-table-example
\textbf{Label: Example: (xy)n = xnyn}
\[(xy)3 = x^{3}y^{3}\]
% formula_record_end

% formula_record_start
% formula_id: exponent-laws-html-027
% index_global: 27
% index_in_page: 27
% section: Laws of Exponents
% subsection: none
% formula_name: (x/y)n = xn/yn
% source: mathsisfun-table-law
\textbf{Label: (x/y)n = xn/yn}
\[(\frac{x}{y})^{n} = xn/yn\]
% formula_record_end

% formula_record_start
% formula_id: exponent-laws-html-028
% index_global: 28
% index_in_page: 28
% section: Laws of Exponents
% subsection: none
% formula_name: Example: (x/y)n = xn/yn
% source: mathsisfun-table-example
\textbf{Label: Example: (x/y)n = xn/yn}
\[(x/y)2 = x^{2} / y^{2}\]
% formula_record_end

% formula_record_start
% formula_id: exponent-laws-html-029
% index_global: 29
% index_in_page: 29
% section: Laws of Exponents
% subsection: none
% formula_name: x-n = 1/xn
% source: mathsisfun-table-law
\textbf{Label: x-n = 1/xn}
\[x^{-n} = 1/xn\]
% formula_record_end

% formula_record_start
% formula_id: exponent-laws-html-030
% index_global: 30
% index_in_page: 30
% section: Laws of Exponents
% subsection: none
% formula_name: Example: x-n = 1/xn
% source: mathsisfun-table-example
\textbf{Label: Example: x-n = 1/xn}
\[x-3 = 1/x^{3}\]
% formula_record_end

% formula_record_start
% formula_id: exponent-laws-html-031
% index_global: 31
% index_in_page: 31
% section: Laws of Exponents
% subsection: none
% formula_name: xm/n = n√xm = (n√x )m
% source: mathsisfun-table-law
\textbf{Label: xm/n = n√xm = (n√x )m}
\[x^{m/n} = n\sqrt{}xm = (n\sqrt{}x )m\]
% formula_record_end

% formula_record_start
% formula_id: exponent-laws-html-032
% index_global: 32
% index_in_page: 32
% section: Laws of Exponents
% subsection: none
% formula_name: Example: xm/n = n√xm = (n√x )m
% source: mathsisfun-table-example
\textbf{Label: Example: xm/n = n√xm = (n√x )m}
\[x\\frac{2}{3} = 3\sqrt{}x^{2} = (3\sqrt{}x )2\]
% formula_record_end

% formula_record_start
% formula_id: exponent-laws-html-033
% index_global: 33
% index_in_page: 33
% section: Uncategorized
% subsection: none
% formula_name: 
% source: formula-text-fallback
\[In this example: 8^{2} = 8 \times 8 = 64\]
% formula_record_end

% formula_record_start
% formula_id: exponent-laws-html-034
% index_global: 34
% index_in_page: 34
% section: Laws Explained
% subsection: none
% formula_name: 
% source: formula-text-fallback
\[The first three laws above (x^{1} = x, x^{0} = 1 and x^{-1} = 1/x) are just part of the natural sequence of exponents. Have a look at this:\]
% formula_record_end

% formula_record_start
% formula_id: exponent-laws-html-035
% index_global: 35
% index_in_page: 35
% section: The law that xmxn = xm+n
% subsection: none
% formula_name: 
% source: formula-text-fallback
\[With x^{m}x^{n}, how many times do we end up multiplying "x"? Answer: first "m" times, then by another "n" times, for a total of "m+n" times.\]
% formula_record_end

% formula_record_start
% formula_id: exponent-laws-html-036
% index_global: 36
% index_in_page: 36
% section: The law that xm/xn = xm-n
% subsection: none
% formula_name: 
% source: formula-text-fallback
\[(Remember that x/x = 1, so every time you see an x "above the line" and one "below the line" you can cancel them out.)\]
% formula_record_end

% formula_record_start
% formula_id: exponent-laws-html-037
% index_global: 37
% index_in_page: 37
% section: The law that xm/xn = xm-n
% subsection: none
% formula_name: 
% source: formula-text-fallback
\[This law can also show you why x^{0}=1 :\]
% formula_record_end

% formula_record_start
% formula_id: exponent-laws-html-038
% index_global: 38
% index_in_page: 38
% section: The law that xm/n = n√xm = (n√x )m
% subsection: none
% formula_name: 
% source: formula-text-fallback
\[x^{1/n} = The n-th Root of x\]
% formula_record_end

% formula_record_start
% formula_id: exponent-laws-html-039
% index_global: 39
% index_in_page: 39
% section: The law that xm/n = n√xm = (n√x )m
% subsection: none
% formula_name: 
% source: formula-text-fallback
\[And so a fractional exponent like 4^{\\frac{3}{2}} is really saying to do a cube (3) and a square root (\\frac{1}{2}), in any order, like this:\]
% formula_record_end

% formula_record_start
% formula_id: exponent-laws-html-040
% index_global: 40
% index_in_page: 40
% section: The law that xm/n = n√xm = (n√x )m
% subsection: none
% formula_name: 
% source: formula-text-fallback
\[Just remember from fractions that m/n = m \times (1/n):\]
% formula_record_end

% formula_record_start
% formula_id: exponent-laws-html-041
% index_global: 41
% index_in_page: 41
% section: The law that xm/n = n√xm = (n√x )m
% subsection: none
% formula_name: 
% source: formula-text-fallback
\[The order does not matter, so it also works for m/n = (1/n) \times m:\]
% formula_record_end

% formula_record_start
% formula_id: exponent-laws-html-042
% index_global: 42
% index_in_page: 42
% section: And That Is It!
% subsection: none
% formula_name: 
% source: formula-text-fallback
\[A fractional exponent like 1/n means to take the nth root: x^{(1n)} = n\sqrt{}x\]
% formula_record_end